
\documentclass[../../../physics-standard_answers_all.tex]{subfiles}

\begin{document}

\setlength{\abovedisplayskip}{2mm}
\setlength{\belowdisplayskip}{1mm}

キルヒホッフ則より，
    %
    \begin{align*}
        V_0 = \frac{Q}{C} + RI \;. \myleader{1.8cm} \text{①}
    \end{align*}
    %
コンデンサの上側極板へ単位時間あたりに流入する電荷が$I$であるから，
    \begin{align*}
        I = \dv{Q}{t} \;. \myleader{1.8cm} \text{②}
    \end{align*}
    %

\begin{enumerate}

    \item $t=0$直後には$Q=0$であるから，①より，
        %
        \begin{align*}
            I_0 = \frac{V_0}{R} \;.
        \end{align*}
        %
        よって，
        %
        \begin{align*}
            R = \frac{V_0}{I_0}
            = \frac{5.0 \, \text{V}}{2.0 \, \text{mA}}
            = \uwave{2.5 \, \text{k}\Omega} \;.
        \end{align*}
        %
        また，充分に時間が経過し，電荷が一定となったときには$I=0$であるから，①より，
        %
        \begin{align*}
            Q_\infty = CV_0 \;.
        \end{align*}
        %
        よって，
        %
        \begin{align*}
            C = \frac{Q_\infty}{V_0}
            = \frac{15 \, \text{\textmu C}}{5.0 \, \text{V}}
            = \uwave{3.0 \, \text{\textmu F}} \;.
        \end{align*}

    \item $Q = kQ_\infty = kCV_0$となったとき，①より，
        %
        \begin{align*}
            V_0 = kV_0 + RI
            \qquad \text{∴} \; I = \uwave{(1-k) \frac{V_0}{R}} \;.
        \end{align*}

    \item ①，②より，
        %
        \begin{align*}
            \dv{Q}{t} = - \frac{1}{RC} (Q-CV_0) \;.
        \end{align*}
        %
        よって，
        %
        \begin{align*}
            Q = \uwave{CV_0 \qty( 1 - e^{-t/RC} )} \;,
            \quad I = \uwave{\frac{V_0}{R} e^{-t/RC}} \;.
        \end{align*}

\end{enumerate}

\end{document}
